\documentclass[10pt,a4paper]{amsart}
\usepackage[margin=19mm]{geometry}
\usepackage{amsmath,amssymb,mathtools}
\usepackage[scr=rsfs,cal=euler]{mathalfa}
\usepackage{xfrac}
\usepackage[unicode,hidelinks,bookmarks=false]{hyperref}
\newcommand{\ie}{i.e.\ }
\newcommand{\cf}{cf.\ }
\newcommand{\resp}{respectively\ }
\newcommand{\Subsection}{\subsection}
\providecommand{\nobreakdash}{\nobreak\hbox{-}}
\setlength{\emergencystretch}{2em}
\multlinegap0pt
\usepackage{amssymb}
\usepackage{mathtools}
\usepackage{xcolor}
\usepackage[shortlabels]{enumitem}
\usepackage{tcolorbox}
\newtheorem{theorem}{Theorem}
\newcommand{\Gspace}{G_\sigma^{p/2}}
\newcommand{\dd}{\mathrm{d}}
\title{\textbf \large {SPATIAL AND TEMPORAL ANALYTICITY OF SOLUTIONS TO SEMI-LINEAR PARABOLIC SYSTEMS WITH ANALYTIC NONLINEARITIES}}
\author{Nubogh B. Qumsiyeh Al-Atrash, Edriss S. Titi}
\date{Source: arXiv:2608.21630v1 (CC BY 4.0)}
\begin{document}
\maketitle

\noindent\emph{Source and licence.} \textbf \large {SPATIAL AND TEMPORAL ANALYTICITY OF SOLUTIONS TO SEMI-LINEAR PARABOLIC SYSTEMS WITH ANALYTIC NONLINEARITIES}. Version arXiv:2608.21630v1 (CC BY 4.0), distributed by arXiv under the Creative Commons Attribution 4.0 International licence, http://creativecommons.org/licenses/by/4.0/, as recorded in the arXiv licence metadata for this exact version and shown on the abstract page https://arxiv.org/abs/2608.21630v1. Full licence notice: \texttt{licenses/arxiv-2608.21630-CC-BY-4.0.txt}.\par\smallskip

\noindent\textit{Editorial scope.} Counted unit: the article's theorem theorem (source0.tex lines 570--580). The article's complete original proof of it is delivered with it (source0.tex lines 582--612, one \texttt{proof} environment of 31 lines). No mathematics is written, repaired or re-derived: the statement and the proof are the article's own lines, in the article's own notation, and no retained line is substituted or re-flowed.  Retained uncounted as the source-local context the proof needs: lines 330--335; lines 371--375; lines 492--495; lines 504--506.  Nothing was omitted from any retained span, so the package carries no omission record.  No retained line contains a citation, so the wrapper carries no bibliography and no bibliography entry.  No retained line contains a figure environment, an image-inclusion command or drawing code, so no image is delivered and the package declares no asset dependency.  Editorial formatting only, in addition: an AMS \texttt{amsart} preamble that restates the article's own declarations and macros used by the retained lines, namely the environments \texttt{theorem} on separate counters, together with the standard packages those lines need.  The retained environments print in this document as Theorem 1 in order of appearance, whereas the article prints the same items with the numbering of its own full document.  The complete native source of the article as distributed in the arXiv e-print for this exact version is bundled unchanged as \texttt{sources/208261/source0.tex}.
\begin{align}
& \frac{1}{2} \frac{\dd}{\dd s} |u^N|_{G_\sigma^{p/2}}^2 + (\nu \cos\theta -\epsilon)|u^N|_{G_\sigma^{\frac{p+1}{2}}}^2 -C_{\epsilon}|u^N|_{
G_\sigma^{p/2}
}^2 \notag \\
&\leq -\operatorname{Re}\left( e^{i\theta}( A^{p/2} e^{\sigma A^{1/2}} P_N F(u^N), A^{p/2} e^{\sigma A^{1/2}} u^N)_ {L^2(\Omega)}\right).\label{11}
\end{align}

\begin{align}
    y(s) &\leq \int_0^s  \!\bigl(1 + C_p^{-1}\bigr)\,
    g \bigl(C_p y(t)\bigl) dt+ \int_0^s  \frac{1}{\nu \cos\theta} y(t)dt
    + y(0).\label{12}
\end{align}

\begin{align}
    \frac{du}{d\zeta} + \nu Au + F(u,\nabla u) &= 0 .\label{20}\\
    u(0) &= u_0,  \label{21}
\end{align}

\begin{align}
    g(r,\rho_1,...,\rho_n)=\sum_\beta |a_\beta| r^{\beta_ 1}\rho_1^{\beta_ 2}...\rho_n^{\beta_{n+1}}, \label{22}
\end{align}

\begin{theorem}
Let $u_0 \in {H}^p_{\text{per}}(\Omega)$ with \( p > \frac{n}{2}+1 \) and $|u_0|_{{H}^p_{\text{per}}(\Omega)} \le M_0$, for some $M_0 > 0$. There exists a continuous, positive function $T^*(\theta)$ defined for $|\theta| < \frac{\pi}{2}$ that depends only on the parameters $M_0$, $\nu$, and $L$, such that equation \eqref{20} with initial condition \eqref{21} has a unique regular solution $u$ inside the domain:
\[
\mathcal{D}_{T^*} := \left\{\zeta = s e^{i\theta} \in \mathbb{C} :\ 0 < s < T^*(\theta), \ \ |\theta| < \frac{\pi}{2}\right\},
\]
and the map
\[
\zeta \longmapsto e^{\sigma A^{1/2}} A^{p/2} u(\zeta), \quad \text{with } \sigma = \sigma(s,\theta) = s\cos\theta,
\]
is analytic in $\mathcal{D}_{T^*}$.
\end{theorem}

\begin{proof}
As the proof strategy mirrors that of Theorems 1 and 2, we present only the key details.
Following Inequality \eqref{11} we have
\begin{align}
& \frac{1}{2} \frac{\dd}{\dd s} |u^N|_{G_\sigma^{p/2}}^2 + (\nu \cos\theta -\epsilon)|u^N|_{G_\sigma^{\frac{p+1}{2}}}^2 -C_{\epsilon}|u^N|_{
G_\sigma^{p/2}
}^2 \notag \\
&\leq -\operatorname{Re}\left( e^{i\theta}
\left( A^\frac{p-1}{2} e^{\sigma A^{1/2}} P_N F(u^N,\nabla u^N), A^\frac{p+1}{2} e^{\sigma A^{1/2}} u^N \right)_ {L^2(\Omega)}\right),\label{26}
\end{align}
where $\epsilon$ to be chosen later.

Using Lemma 1, Lemma 2, Young's inequality, the fact that  $|\nabla u|^2_{G_{\sigma}^{\frac{p-1}{2}}}\leq | u|^2_{G_{\sigma}^{\frac{p}{2}}}$  and the majorizing function $g$ as in \eqref{22}, the absolute value of the  right-hand side of \eqref{26}is
\begin{align}
&\Big|-\operatorname{Re}\left( e^{i\theta}
\left( A^{\frac{p-1}{2}} e^{\sigma A^{1/2}} P_N F(u^N,\nabla u^N), A^{\frac{p+1}{2}} e^{\sigma A^{1/2}} u^N \right)_ {L^2(\Omega)}\right)\Big| \notag \\
&\quad\leq  |P_NF(u^N,\nabla u^N)|_{G^{\frac{p-1}{2}}_{\sigma}}|u^N|_{G^{\frac{p+1}{2}}_{\sigma}}\notag \\
&\quad \leq  |F(u^N,\nabla u^N)|_{G^{\frac{p-1}{2}}_{\sigma}}|u^N|_{G^{\frac{p+1}{2}}_{\sigma}}\notag \\
&\quad \leq C_{\epsilon}|F(u^N,\nabla u^N)|^2_{G^{\frac{p-1}{2}}_{\sigma}}+\epsilon|u^N|^2_{G^{\frac{p+1}{2}}_{\sigma}}\notag \\
&\quad \leq C_{\epsilon}C g^2(C_1 |u^N|_{G_\sigma^{\frac{p-1}{2}}},C_2 |u^N_{x_1}|_{G_\sigma^{\frac{p-1}{2}}},...,C_{n+1} |u^N_{x_n}|_{G_\sigma^{\frac{p-1}{2}}})+\epsilon|u^N|^2_{G^{\frac{p+1}{2}}_{\sigma}}\notag \\
&\quad \leq C_{\epsilon}C g^2(C_1 |u^N|_{G_\sigma^{\frac{p}{2}}},C_2 |u^N|_{G_\sigma^{\frac{p}{2}}},...,C_{n+1} |u^N|_{G_\sigma^{\frac{p}{2}}})+\epsilon|u^N|^2_{G^{\frac{p+1}{2}}_{\sigma}}.\label{27}
\end{align}
Choosing $\epsilon=\dfrac{\nu \cos\theta}{4},C_{\epsilon}=\dfrac{1}{\nu \cos\theta}$, from \eqref{26} and \eqref{27}, we obtain:
\begin{align*}
\frac{1}{2} \frac{\dd}{\dd s} |u^N|^2_{\Gspace} + \frac{\nu\cos\theta}{2}|u^N|_{G_\sigma^{\frac{p+1}{2}}}^2 &\leq \frac{C}{\nu \cos\theta} g^2(C_1 |u^N|_{G_\sigma^{\frac{p}{2}}},C_2 |u^N|_{G_\sigma^{\frac{p}{2}}},...,C_{n+1} |u^N|_{G_\sigma^{\frac{p}{2}}})\notag\\&\quad+\frac{1}{\nu \cos\theta}|u^N|_{
G_\sigma^{p/2}
}^2.
\end{align*}

Now, we proceed as in the proof of Theorem 1, starting from \eqref{12} to finish the proof of Theorem 3.
\end{proof}

\end{document}
